\SourceLeaf{259}%
\(DB\). For this point, \(p\) remains the same, and the symbols \(q',r',\varphi',\psi',S'\) have the meanings that \(q,r,\varphi,\psi,S\) have for the point \(B\). This gives a triangle with vertices \(A,B,C\). We denote its angles by \(A,B,C\), the sides opposite those angles by \(a,b,c\), and the area by \(\sigma\). We represent the measure of curvature at \(A,B,C\) by \(\alpha,\beta,\gamma\), respectively. If we assume, as we may, that \(p,q,q-q'\) are positive, we have
\[
\begin{gathered}
A=\varphi-\varphi',\quad B=\psi,\quad C=\pi-\psi',\\
a=q-q',\quad b=r',\quad c=r,\quad\sigma=S-S'
\end{gathered}
\]

First we shall express the area \(\sigma\) as a series. Replacing each quantity in [7] that refers to \(B\) by its counterpart for \(C\) gives the formula for \(S'\). From this we obtain, up to quantities of the sixth order
\[
\begin{aligned}
\sigma=\tfrac{1}{2}p(q-q')\{&1-\tfrac{1}{6}f^0(pp+qq+qq'+q'q')\\
&-\tfrac{1}{60}f'p(6pp+7qq+7qq'+7q'q')\\
&-\tfrac{1}{20}g^0(q+q')(3pp+4qq+4qq'+4q'q')\}
\end{aligned}
\]
With the help of series [2], namely
\[
c\sin B=p(1-\tfrac{1}{3}f^0qq-\tfrac{1}{4}f'pqq-\tfrac{1}{2}g^0q^3-\text{etc.})
\]
this formula takes the following form
\[
\begin{aligned}
\sigma=\tfrac{1}{2}ac\sin B\{&1-\tfrac{1}{6}f^0(pp-qq+qq'+q'q')\\
&-\tfrac{1}{60}f'p(6pp-8qq+7qq'+7q'q')\\
&-\tfrac{1}{20}g^0(3ppq+3ppq'-6q^3+4qqq'+4qq'q'+4q'^3)\}
\end{aligned}
\]

The measure of curvature at any point of the surface becomes (by Article 19, in which \(m,p,q\) stood for what we now denote by \(n,q,p\))
\[
\begin{aligned}
&=-\frac{1}{n}\mathord{.}\frac{\mathrm{d}\mathrm{d}n}{\mathrm{d}q^2}=-\frac{2f+6gq+12hqq+\text{etc.}}{1+fqq+\text{etc.}}\\
&=-2f-6gq-(12h-2ff)qq-\text{etc.}
\end{aligned}
\]
This gives, when \(p,q\) refer to the point \(B\),
\[
\beta=-2f^0-2f'p-6g^0q-2f''pp-6g'pq-(12h^0-2f^0f^0)qq-\text{etc.}
\]
and also
\[
\begin{aligned}
\gamma&=-2f^0-2f'p-6g^0q'-2f''pp-6g'pq'-(12h^0-2f^0f^0)q'q'-\text{etc.}\\
\alpha&=-2f^0
\end{aligned}
\]
